\section{重回帰分析}
\label{lesson09}

\subsection{授業内容}
\subsubsection{科目の中でのこのコマの位置づけ}
\begin{tabular}[t]{ccccccc}
    記述統計[3] & 確率[2] & モデル[3] & 推測・推定[1] & 意思決定[0] & 実践[1]
\end{tabular}

\subsubsection{概要}
回帰分析によって，データにモデルを当てはめて係数の推定値を得るということを行った。データに最も適した形を求めて計算を行ったが，そのこととモデルが適していたかどうかは別である。この点を理解するためには，モデルの適合度について考える必要がある。モデルが適合しているということを踏まえた上で，説明する変数が増える重回帰分析へと議論は展開する。ここでは説明変数の読み取り方への注意を促すべく，偏相関係数の考え方を解説し，偏回帰係数，標準化偏回帰係数などの用語について理解する。
\subsubsection{コマ主題細目(3-5項目箇条書き)}
\begin{description}
\item[残差と決定係数]回帰係数が計算されれば，実際のデータに対して予測値を計算することができ，予測がどの程度適合していたかを検証することができる。予測値と被説明変数との差分を残差residualsと呼ぶが，この両者の相関係数を重相関係数，その二乗したものを決定係数と呼んでモデルの適合度を測る指標にすることを理解する。
\begin{flushright}
    $\to$\citet{野島201903}のP.58-60.
\end{flushright}

\item[重回帰分析]説明変数が複数に増えた場合の回帰分析は，重回帰分析Multiple Regression Analysisと呼ぶ。数式的な表現としては項が増えるだけであるが，データのプロットを考えると三次元以上の回帰平面に拡張されていること，それでも面は湾曲していない，すなわち線形関係であることを理解する。
\begin{flushright}
    $\to$ \citet{kosugi2019}のP.168,図10.7の回帰平面図を参照
\end{flushright}

\item[偏回帰係数]単回帰分析のときとは違って複数の説明変数があることから，単一の回帰係数の解釈は条件付きで考えなければならない。そのことを理解するために，部分相関，偏相関の考え方を理解し，その上で偏回帰係数の読み取り方の理解へとつなげる。また変数を増やすことで重相関係数の増加することを理解する。
\begin{flushright}
    $\to$\citet{kosugi2018}のP.60--63.
\end{flushright}

\item[標準化された係数]回帰係数については，素点による回帰係数が理解しやすいが，単位に左右されるものでもあるので相対比較をする場合は全ての変数を標準化した標準化解を用いた方がわかりやすい。標準化解を用いても適合度は変化しないことなど留意点とともに有用性を理解する。
\begin{flushright}
    $\to$\citet{kosugi2018}のP.66--67.
\end{flushright}


\end{description}
\subsubsection{キーワード}
\begin{itemize}
\item 残差
\item 決定係数
\item 重回帰分析
\item 偏相関係数，偏回帰係数
\item 標準化解
\end{itemize}
\subsection{授業情報}
\paragraph{コマの展開方法}
講義
\paragraph{標準シラバスにおける位置づけ}
科目番号5 心理学統計法；1.心理学で用いられる統計手法；C 記述統計：回帰
\subsubsection{予習・復習課題}
\paragraph{予習}
説明変数が一つの回帰分析について，数式的表現や最適化関数(最小二乗法)の意味について概念的な理解で十分なので再確認しておくこと。後者について今回も同じ基準が用いられるが，説明変数が複数になる場合は解析的算出には行列計算の知識が必要である。行列計算については，二年次の統計に関する講義で扱う予定であるが，興味があるものは「線形代数」をキーワードに予習すると良い。初学者には\citet{結城201810}が，あるいは\citet{岡太200804}がよい。
\paragraph{復習}
\citet{豊田201706}のP.89--98は本時の内容がわかりやすくまとめられているので，通読しておくと良い復習になる。また重回帰分析は広く用いられる手法なので，\citet{kosugi2016a}の第二章などを読んで，どのような使い方をされているのか，また前提となっている条件や利用上の問題点などについて，講義時間内で触れられなかった諸問題についての知識を補完しておいてもらいたい。
